\documentclass[11pt]{article}
\usepackage[letterpaper,margin=0.85in]{geometry}\usepackage[T1]{fontenc}\usepackage[utf8]{inputenc}\usepackage{amsmath,amssymb,enumitem,booktabs,array,xcolor,fancyhdr,tabularx}
\setlist{itemsep=2pt,topsep=3pt}\pagestyle{fancy}\fancyhf{}\renewcommand{\headrulewidth}{0.3pt}
\fancyhead[L]{\small 4.3 Rates of Change in Other Contexts}\fancyhead[R]{\small Make-up key}\fancyfoot[C]{\small\thepage}
\definecolor{ans}{RGB}{28,125,74}\newcommand{\A}[1]{\textcolor{ans}{\textbf{#1}}}
\begin{document}

\section*{4.3 Rates of Change in Other Contexts}
\textbf{You missed this lesson. Copy this into your notes. That is all you need to do.}\par\medskip
\subsection*{Question of the day}
How do average rates and instant rates work for things that are not moving objects?\par
\subsection*{Targets}\begin{itemize}\item I can find an average rate of change from a table or a formula.
\item I can tell an average rate from a rate at one moment.
\item I can use f? and f? to say how a quantity is changing.\end{itemize}
\subsection*{Concepts}\begin{itemize}
\item \textbf{Average rate of change} ((f(b) -- f(a)) / (b -- a)): Change in f divided by change in x over an interval. It is the slope of the secant line through the two ends.
\item \textbf{Rate at one moment} (f?(a)): The derivative f?(a) is the rate at one instant. It is the slope of the tangent line at x = a.
\item \textbf{Units} (units of f per unit of x): Both kinds of rate have units of f per unit of x. Say "per hour", "per year" and so on.
\item \textbf{The rate of the rate} (f? = (f?)?): f? says whether the rate f? is getting bigger or smaller. f? > 0 and f? < 0: still rising, but more and more slowly.
\end{itemize}
\subsection*{What the class predicted, and what is true}\begin{itemize}
\item A savings account holds $200 after 2 years and $260 after 5 years. On average, it grew by... \A{20 dollars per year.} Change over time: $\dfrac{260-200}{5-2}=20$ dollars per year.
\item Is the average rate of change over an interval always the same as the rate at one moment inside it? \A{No: the rate can change inside the interval.} On a straight line they match. On a curve the rate changes, so an average can differ from the rate at one moment.
\end{itemize}
\section*{Huddle notes (answers filled in)}
\subsection*{Definitions}\begin{itemize}
\item \textbf{Average rate of change:} From x = a to x = b: \A{change in f} divided by \A{change in x}. It is the slope of the secant line.
\item \textbf{Rate at one moment:} f?(a) is the rate at one \A{instant}. It is the slope of the \A{tangent} line.
\item \textbf{Units:} Both rates have units of f \A{per} unit of x, like "thousand bacteria per hour".
\item \textbf{The rate of the rate:} If f? > 0 the rate is getting \A{bigger}. If f? < 0 the rate is getting \A{smaller}.
\end{itemize}
\subsection*{Worked example}
N(t) = --t\textsuperscript{2} + 12t + 20 thousand bacteria, t in hours. Find the average rate from t = 2 to t = 6, and the rate at t = 2.\\
\begin{itemize}
\item N(2) = 40 and N(6) = \A{56}
\item Average rate = (56 -- 40) / 4 = \A{4} thousand per hour
\item N?(t) = --2t + \A{12}
\item Rate at t = 2: N?(2) = \A{8} thousand per hour
\end{itemize}
\subsection*{Label the picture}
\begin{enumerate}
\item The red tangent line: its slope is the \A{instantaneous} rate.
\item The orange secant line: its slope is the \A{average} rate.
\item The two points mark the \A{ends} of the interval.
\end{enumerate}
\section*{One worked sample of each skill}
\subsection*{Average rate of change}
\noindent\textbf{Level 1.} $P(t)$ is the population of a city, in thousands, $t$ years after 2000. You know $P(3)=260$ and $P(8)=280$. Find the average rate of change of $P$ from $t=3$ to $t=8$.\\ \A{Answer: 4 thousand people per year} \quad $\dfrac{280-260}{8-3}=\dfrac{20}{5}=4$ thousand people per year.\par\medskip
\noindent\textbf{Level 2.} $N(t)$ is the number of bacteria, in thousands, in a dish $t$ hours after noon. The table gives some values. Find the average rate of change of $N$ from $t=3$ to $t=7$.\\ \A{Answer: 3 thousand bacteria per hour} \quad $\dfrac{37-25}{7-3}=\dfrac{12}{4}=3$.\par\medskip
\noindent\textbf{Level 3.} $S(d)=d^{3}-2d^{2}+11d+10$ is the depth, in inches, of the snow on the ground $d$ days after a storm. Find the average rate of change of $S$ from $d=2$ to $d=4$.\\ \A{Answer: 27 inches per day} \quad $S(2)=32$ and $S(4)=86$. So $\dfrac{86-32}{2}=27$.\par\medskip
\noindent\textbf{Level 4.} $S(d)=d^{2}+3d+56$ is the depth, in inches, of the snow on the ground $d$ days after a storm. Find the average rate of change of $S$ from $d=2$ to $d=4$.\\ \A{Answer: 9 inches per day} \quad $S(2)=66$ and $S(4)=84$. So $\dfrac{84-66}{2}=9$.\par\medskip
\subsection*{Rate at one moment}
\noindent\textbf{Level 1.} A cube has side $s$ inches. Its volume is $V(s)=s^3$ cubic inches. Find $V'(4)$: how fast does it change per inch of side when the side is 4 inches?\\ \A{Answer: 48 cubic inches per inch} \quad $V'(s)=3s^2$. Putting in 4 gives $48$ cubic inches per inch.\par\medskip
\noindent\textbf{Level 2.} A circle has radius $r$ inches. Its area is $A(r)=\pi r^2$ square inches. Find $A'(6)$ and give your answer to 2 decimal places.\\ \A{Answer: 37.7 square inches per inch} \quad $A'(r)=2\pi r$. At 6 inches this is about $37.7$ square inches per inch.\par\medskip
\noindent\textbf{Level 3.} The amount of a medicine in the blood is $C(t)=100e^{-0.1t}$ milligrams after $t$ hours. Find $C'(0)$ to 2 decimal places. The answer is negative: the quantity is falling.\\ \A{Answer: -10 milligrams per hour} \quad $C'(t)=-10e^{-0.1t}$. At that time it is about $-10$ milligrams per hour.\par\medskip
\noindent\textbf{Level 4.} The amount of a medicine in the blood is $C(t)=40e^{-0.3t}$ milligrams after $t$ hours. Find $C'(10)$ to 2 decimal places. The answer is negative: the quantity is falling.\\ \A{Answer: -0.597 milligrams per hour} \quad $C'(t)=-12e^{-0.3t}$. At that time it is about $-0.597$ milligrams per hour.\par\medskip
\subsection*{How is it changing?}
\noindent\textbf{Level 1.} $P(t)$ is the population of a city, in thousands, $t$ years after 2000. At $t=7$ we know $P'(7)=-4$ and $P''(7)=2$. What is happening to the population?\\ (A) increasing, and the rate is getting bigger \quad (B) decreasing, and it is falling faster and faster \quad \A{(C) decreasing, but it is falling more and more slowly} \quad (D) increasing, but the rate is getting smaller\\ \A{Answer: C} \quad $P'$ is negative, so the population is decreasing. $P''$ is positive, so the rate is getting bigger.\par\medskip
\noindent\textbf{Level 2.} $N(t)=-t^{3}-4t^{2}+9t+60$ is the number of bacteria, in thousands, in a dish $t$ hours after noon. At $t=1$, what is happening to the number of bacteria?\\ (A) decreasing, but it is falling more and more slowly \quad (B) increasing, but the rate is getting smaller \quad (C) increasing, and the rate is getting bigger \quad \A{(D) decreasing, and it is falling faster and faster}\\ \A{Answer: D} \quad $N'(1)=-2$ and $N''(1)=-14$. Decreasing, and it is falling faster and faster.\par\medskip
\noindent\textbf{Level 3.} $S(d)$ is the depth, in inches, of the snow on the ground $d$ days after a storm. The rate $S'$ is -9 at $d=1$ and -5 at $d=5$ (inches per day), and it changes steadily in between. At $d=1$, what is happening to the snow depth?\\ (A) increasing, and the rate is getting bigger \quad (B) decreasing, and it is falling faster and faster \quad \A{(C) decreasing, but it is falling more and more slowly} \quad (D) increasing, but the rate is getting smaller\\ \A{Answer: C} \quad The sign of $S'$ is negative. The rate goes from -9 to -5, so it is getting bigger ($S''>0$). So decreasing, but it is falling more and more slowly.\par\medskip
\noindent\textbf{Level 4.} $S(d)$ is the depth, in inches, of the snow on the ground $d$ days after a storm. The rate $S'$ is 10 at $d=4$ and 13 at $d=6$ (inches per day), and it changes steadily in between. At $d=4$, what is happening to the snow depth?\\ (A) decreasing, but it is falling more and more slowly \quad (B) increasing, but the rate is getting smaller \quad \A{(C) increasing, and the rate is getting bigger} \quad (D) decreasing, and it is falling faster and faster\\ \A{Answer: C} \quad The sign of $S'$ is positive. The rate goes from 10 to 13, so it is getting bigger ($S''>0$). So increasing, and the rate is getting bigger.\par\medskip
\subsection*{Say what the average means}
\noindent\textbf{Level 1.} $P(t)$ is the population of a city, in thousands, $t$ years after 2000. The average rate of change of $P$ from $t=2$ to $t=6$ is $8$ thousand people per year. What does this say?\\ (A) At exactly t = 2, the population is rising at 8 thousand people per year. \quad (B) At t = 6, the population is 8 thousand people. \quad \A{(C) On average, the population grew by 8 thousand people per year from t = 2 to t = 6.} \quad (D) The population grew by a total of 8 thousand people from t = 2 to t = 6.\\ \A{Answer: C} \quad The sentence must name the interval (2 to 6), say "on average", and give the rate per year: On average, the population grew by 8 thousand people per year from t = 2 to t = 6.\par\medskip
\noindent\textbf{Level 2.} $N(t)$ is the number of bacteria, in thousands, in a dish $t$ hours after noon. The average rate of change of $N$ from $t=5$ to $t=9$ is $-2$ thousand bacteria per hour. What does this say?\\ (A) At exactly t = 5, the number of bacteria is falling at 2 thousand bacteria per hour. \quad (B) The number of bacteria fell by a total of 2 thousand bacteria from t = 5 to t = 9. \quad (C) At t = 9, the number of bacteria is 2 thousand bacteria. \quad \A{(D) On average, the number of bacteria fell by 2 thousand bacteria per hour from t = 5 to t = 9.}\\ \A{Answer: D} \quad The sentence must name the interval (5 to 9), say "on average", and give the rate per hour: On average, the number of bacteria fell by 2 thousand bacteria per hour from t = 5 to t = 9.\par\medskip
\noindent\textbf{Level 3.} $S(d)$ is the depth, in inches, of the snow on the ground $d$ days after a storm. Which equation says: "the average rate of change of $S$ from $d=1$ to $d=5$ is 2 inches per day"?\\ (A) $S'(1)=2$ \quad (B) $\dfrac{S(5)-S(1)}{1}=2$ \quad \A{(C) $\dfrac{S(5)-S(1)}{5-1}=2$} \quad (D) $S(5)-S(1)=2$\\ \A{Answer: C} \quad Average rate = change in $S$ $\div$ change in $d$ = $\dfrac{S(5)-S(1)}{5-1}$.\par\medskip
\noindent\textbf{Level 4.} $S(d)$ is the depth, in inches, of the snow on the ground $d$ days after a storm. Which equation says: "the average rate of change of $S$ from $d=5$ to $d=7$ is 8 inches per day"?\\ (A) $\dfrac{S(7)-S(5)}{5}=8$ \quad (B) $S'(5)=8$ \quad \A{(C) $\dfrac{S(7)-S(5)}{7-5}=8$} \quad (D) $S(7)-S(5)=8$\\ \A{Answer: C} \quad Average rate = change in $S$ $\div$ change in $d$ = $\dfrac{S(7)-S(5)}{7-5}$.\par\medskip
\end{document}